\subsection{The approximation theorem for Dedekind domains}

In Dedekind domains there is an ``approximation theorem'' which strengthens both Theorem 1 of Chapter VI, § 7, no. 2 and Proposition 9 of § 1, no. 5:

PROPOSITION 2. Let A be a Dedekind domain, K its field of fractions and P the set of non-zero prime ideals of A; for \(\mathfrak{p} \in \mathbf{P}\) let \(v_{\mathfrak{p}}\) denote the corresponding essential valuation of A. Let \(\mathfrak{p}_1, \ldots, \mathfrak{p}_q\) be distinct elements of P and \(n_1, \ldots, n_q\) rational integers and \(x_1, \ldots, x_q\) elements of K. Then there exists \(x \in \mathbf{K}\) such that \(v_{\mathfrak{p}_i}(x - x_i) \geqslant n_i\) for \(1 \leqslant i \leqslant q\) and \(v_{\mathfrak{p}}(x) \geqslant 0\) for all \(\mathfrak{p} \in \mathbf{P}\) distinct from the \(\mathfrak{p}_i\) .

Replacing if need be the \(n_{i}\) by greater integers, they may be assumed all to be positive. We examine first the case where the \(x_{i}\) are in A; it obviously amounts to finding an \(x \in A\) satisfying the congruences

\[
x \equiv x _ {i} (\mathrm{mod.} p _ {i} ^ {n _ {i}})
\]

and the existence of \(x\) then follows from Chapter II, § 1, no. 2, Proposition 5.

We pass now to the general case. We may write \(x_{i} = s^{-1}y_{i}\) where \(s, y_{i}\) are in A; writing \(x = s^{-1}y\) , it amounts to finding a \(y \in A\) such that, on the one hand, \(v_{p_i}(y - y_i) \geqslant n_i + v_{p_i}(s)\) and, on the other, \(v_{p}(y) \geqslant v_{p}(s)\) for all \(p \in P\) distinct from the \(p_i\) ; as \(v_{p}(s) = 0\) except for a finite number of indices \(p\) , it is thus reduced to the above case; whence the proposition.

Proposition 2 may be interpreted as a density theorem. To be precise, for all \(p \in P\) , let \(\hat{K}_p\) (resp. \(\hat{A}_p\) ) be the completion of \(K\) (resp. \(A\) ) with respect to the discrete valuation \(v_p\) and consider the product \(\prod_{p \in P} \hat{K}_p\) ; an element \(x = (x_p)\) of this product is called a restricted adèle of \(A\) if \(x_p \in \hat{A}_p\) for all \(p \in P\) with the exception of a finite number of them. Clearly the set \(A\) of restricted adèles is a subring of \(\prod_{p \in P} \hat{K}_p\) , which contains the product ring \(A_0 = \prod_{p \in P} \hat{A}_p\) . Consider on \(A_0\) the product topology, with respect to which \(A_0\) is complete; there is on \(A\) a unique topology \(T\) which is compatible with its additive group structure and for which the neighbourhoods of 0 in \(A_0\) form a fundamental system \(S\) of neighbourhoods of 0. The topology \(T\) is compatible with the ring structure on \(A\) ; for clearly axiom ( \(AV_{II}\) ) of General Topology, Chapter III, § 6, no. 3, holds, the topology induced by \(T\) on \(A_0\) being compatible with the ring structure on \(A_0\) . On the other hand, for all \(x \in A\) there exists a finite subset \(J\) of \(P\) such that, if we write \(J' = P - J\) , \(K_J \leqslant \prod_{p \in J} \hat{K}_p\) , \(A_{J'} = \prod_{p \in J'} \hat{A}_p\) , then \(x \in K_{J} \times A_{J'}\) and, as \(\hat{A}_{p}\) is open in \(\hat{K}_{p}\) for all p, \(\mathfrak{S}\) is a fundamental system of neighbourhoods of 0 for the product topology on \(K_{J} \times A_{J'}\) ; since the latter is compatible with the ring structure on this product, axiom (AV \(_{I}\) ) of General Topology, Chapter III, loc. cit. is also seen to hold, which proves our assertion. Clearly \(A_{0}\) is an open subring of A and hence A is also a complete ring (General Topology, Chapter III, § 3, no. 3, Proposition 4).

For all \(x \in \mathbf{K}\) , let \(\Delta(x)\) be the element \((x_{\mathfrak{p}}) \in \bigcup_{\mathfrak{p} \in \mathbb{P}} \hat{\mathbf{K}}_{\mathfrak{p}}\) such that \(x_{\mathfrak{p}} = x\) for all \(\mathfrak{p} \in \mathbb{P}\) ; as \(x_{\mathfrak{p}} \in \hat{\mathbf{A}}_{\mathfrak{p}}\) except for a finite number of values of \(\mathfrak{p}, \Delta(x) \in A\) ; hence we have thus defined a homomorphism \(\Delta: \mathbf{K} \to A\) which is injective if \(\mathbf{P} \neq \varnothing\) (that is if \(A\) is not a field); the elements of \(\Delta(\mathbf{K})\) are called principal restricted adèles and clearly \(\Delta(A) \subset A_0\) .

PROPOSITION 3. The ring \(A_0\) (resp. \(A\) ) is identified with the completion of \(A\) (resp. \(K\) ) with respect to the ring topology for which a fundamental system of neighbourhoods of 0 consists of all the integral ideals \(\neq (0)\) of \(A\) .

It is immediate that the topology considered on A (or K) is Hausdorff. Taking account of no. 3, the assertion concerning \(A_{0}\) follows from Chapter III, § 2, no. 13, Proposition 17. This shows therefore that \(\Delta(A)\) is dense in \(A_{0}\) ; to see similarly that \(\Delta(K)\) is dense in \(A\) , note that for all \(x = (x_{p}) \in A\) there is only a finite number of \(p \in P\) such that \(v_{p}(x_{p}) < 0\) ; by § 1, no. 5, Proposition 9 there is therefore an \(s \in K\) such that \(sx_{p} \in \hat{A}_{p}\) for all \(p \in P\) , in other words \(\Delta(s)x \in A_{0}\) and, as multiplication by \(\Delta(s)\) is a homomorphism from \(A\) to itself, it suffices to apply the fact that \(\Delta(A)\) is dense in \(A_{0}\) to deduce that \(\Delta(K)\) is dense in \(A\) .

We could of course also prove that \(\Delta(\mathbf{K})\) is dense in A by using Proposition 2.

Consider now the multiplicative group \(\mathbf{SL}(n, A)\) consisting of the matrices \(U \in \mathbf{M}_n(A)\) such that \(\det(U) = 1\) ; if \(\mathbf{M}_n(A) = A^{n^2}\) is given the product topology, it induces on \(\mathbf{SL}(n, A)\) a topology compatible with the group structure on \(\mathbf{SL}(n, A)\) . It suffices to verify that the mapping \(U \mapsto U^{-1}\) is continuous on \(\mathbf{SL}(n, A)\) ; but as \(U\) is unimodular, it is known (Algebra, Chapter III, § 6, no. 5, formula (17)) that the elements of \(U^{-1}\) are minors of \(U\) and hence polynomials in the elements of \(U\) , which proves our assertion. If \(K\) is identified with a subring of \(A\) by means of \(\Delta\) , the group \(\mathbf{SL}(n, K)\) is a subgroup of \(\mathbf{SL}(n, A)\) .

PROPOSITION 4. The group \(\mathbf{SL}(n, \mathbf{K})\) is dense in \(\mathbf{SL}(n, A)\) .

Let \(G\) be the closure of \(\mathbf{SL}(n, K)\) in \(\mathbf{SL}(n, A)\) ; as \(K\) is dense in \(A\) (Proposition 3), \(G\) contains all the matrices of the form \(I + a.E_{ij}\) for \(i \neq j\) and \(a \in A\) . For all \(p \in P\) and all \(\lambda \in \hat{K}_p\) , let \(\lambda(p)\) be the restricted adèle \(x = (x_q)_{q \in P}\) such that \(x_p = \lambda\) and \(x_q = 0\) for \(q \neq p\) ; the above shows that \(G\) contains the matrices \(I + \lambda(p)E_{ij}\) for \(i \neq j\) . But we know that the matrices of the form \(I + \lambda E_{ij}\) for \(\lambda \in \hat{K}_p\) generate the group \(\mathbf{SL}(n, \hat{\mathbf{K}}_p)\) (Algebra, Chapter III). For every matrix \(U \in \mathbf{SL}(n, A)\) let \(U_{p}\) denote the canonical image of U in \(\mathbf{SL}(n, \hat{\mathbf{K}}_{\mathfrak{p}})\) ; it is therefore seen that, for all \(p \in P\) , G contains the matrices \(U \in \mathbf{SL}(n, A)\) such that \(U_{q} = I\) for all \(q \neq p\) . Since G is a group, it also contains all the matrices \(U \in \mathbf{SL}(n, A)\) such that \(U_{p} = I\) except for a finite number of \(p \in P\) ; now, the definition of the topology on A shows immediately that the set of these matrices is dense in \(\mathbf{SL}(n, A)\) .

